\section{Introduction}
\IEEEPARstart{L}{ow} Earth orbit (LEO) constellations such as Starlink and OneWeb now bring broadband to aircraft, ships, and remote sites, and 3GPP integrates such non-terrestrial links into 5G~\cite{lin20215g}. A terminal that has a LEO path and a terrestrial path can protect delay-sensitive traffic by sending each packet on the path that will deliver it in time, or on both paths when neither is safe. Multipath transport protocols such as Multipath QUIC (MPQUIC) provide the mechanism~\cite{deconinck2017multipath}. Their schedulers, however, act on smoothed round-trip time (RTT) estimates of the recent past~\cite{paasch2014experimental,lim2017ecf}. On LEO paths this is too slow. Starlink latency jumps at globally synchronized reconfiguration instants every 15\,s~\cite{mohan2024multifaceted,tanveer2023making}. OneWeb latency changes at satellite handovers~\cite{zhao2025measuring}, and measured multipath sessions over Starlink and cellular paths show scheduler failures caused by these dynamics~\cite{baltaci2023multipath}. Schedulers that act on forecasts of path state have therefore been proposed~\cite{yang2024mams,maglione2025intelligent}. Such a scheduler needs two things from a forecaster: an estimate of the delay of each path a few hundred milliseconds to a few seconds ahead, and a bound on how wrong that estimate may be. A late packet costs a missed deadline, and a needless duplicate costs capacity.

Work on LEO performance prediction forecasts throughput or latency at coarse time scales and judges the forecasts by their point accuracy~\cite{liu2025vivisecting,benghe2026horizon}. At the sub-second scale of a packet scheduler this yardstick is discouraging. On cellular links, learned models gain little over simple predictors at fine time scales~\cite{basit2024predictability}. In our data, no model we trained, deep or linear, lowers the RTT root-mean-square error (RMSE) of our previous recurrent forecaster by more than \MaxRMSEGain\% on any link. If point accuracy were the goal, learning would add little. We argue that it is the wrong goal for scheduling and ask a different question. \emph{What} about future link state can be predicted, and how can a scheduler use it?

The answer starts from the data. Fig.~\ref{fig:trace} shows one minute of Starlink downlink RTT forecast 1\,s ahead. The link first stays within a few milliseconds of its level and then starts to swing by about 150\,ms once per second. No forecaster can say where in the swing the RTT will be 1\,s later, but the state of the link is visible in the last seconds of input. A point forecast with an error interval of constant width, the usual way to state uncertainty, must be as wide as the swings at all times. It wastes most of its width while the link is quiet. A forecaster of the distribution of the next RTT can use a narrow interval in the quiet state and a wide one in the loaded state, and it can place most of the width above its median, where RTT spikes go. Across six links we find this picture to be general. The level of the next RTT is hard to predict beyond a return toward its recent value, but its spread differs strongly between windows and can be read from the input.

\begin{figure}[t]
  \centering
  % fr_trace.tex - generated by modeling/experiments/fr_trace.py, do not edit by hand.
\begin{tikzpicture}
\begin{axis}[frbase, name=main, width=8.7cm, height=5.6cm, xmin=0, xmax=59.8, ymin=-50, ymax=300,
  xlabel={time (s)}, ylabel={RTT 1\,s ahead (ms)}, legend cell align=left, legend style={at={(0.5,1.02)}, anchor=south, legend columns=2, font=\scriptsize, draw=none}]
\addplot[name path=sh, draw=none, forget plot] table[x=t, y=prevhi] {figs/data/fr_trace_starlink_dl.dat};
\addplot[name path=sl, draw=none, forget plot] table[x=t, y=prevlo] {figs/data/fr_trace_starlink_dl.dat};
\addplot[cPrev!18] fill between[of=sh and sl]; \addlegendentry{S2S-E constant 95\% band}
\addplot[name path=dh, draw=none, forget plot] table[x=t, y=frhi] {figs/data/fr_trace_starlink_dl.dat};
\addplot[name path=dl, draw=none, forget plot] table[x=t, y=frlo] {figs/data/fr_trace_starlink_dl.dat};
\addplot[cFR!45] fill between[of=dh and dl]; \addlegendentry{DS2S-E calibrated 95\% band}
\addplot[black, thin] table[x=t, y=truth] {figs/data/fr_trace_starlink_dl.dat}; \addlegendentry{measured RTT}
\addplot[cFR, thick] table[x=t, y=frmed] {figs/data/fr_trace_starlink_dl.dat}; \addlegendentry{DS2S-E median}
\draw[black!70] (axis cs:2.0,22.56) rectangle (axis cs:10.0,51.59);
\coordinate (zt) at (axis cs:6,51.59);
\end{axis}
\begin{axis}[at={($(main.north west)+(0.75cm,-0.25cm)$)}, anchor=north west, width=3.6cm, height=2.6cm,
  xmin=2.0, xmax=10.0, ymin=22.56, ymax=51.59, tick label style={font=\tiny}, axis background/.style={fill=white},
  grid=major, grid style={gray!20}, name=ins]
\addplot[name path=dh2, draw=none, forget plot] table[x=t, y=frhi] {figs/data/fr_trace_starlink_dl.dat};
\addplot[name path=dl2, draw=none, forget plot] table[x=t, y=frlo] {figs/data/fr_trace_starlink_dl.dat};
\addplot[cFR!45] fill between[of=dh2 and dl2, soft clip={domain=2:10}];
\addplot[black, thin, restrict x to domain=2.0:10.0] table[x=t, y=truth] {figs/data/fr_trace_starlink_dl.dat};
\addplot[cFR, thick, restrict x to domain=2.0:10.0] table[x=t, y=frmed] {figs/data/fr_trace_starlink_dl.dat};
\end{axis}
\draw[->, black!70] (zt) -- (ins.south east);
\end{tikzpicture}

  \caption{One minute of a Starlink DL test session, forecast 1\,s ahead. The previous model (S2S-E) has a constant 95\% interval of \TraceWidS\,ms. The calibrated interval of the proposed model (DS2S-E) is \TraceWidDCalm\,ms wide (median) while the link is quiet and \TraceWidDLoad\,ms once it swings. Inset: the quiet part.}
  \label{fig:trace}
\end{figure}

This paper turns that observation into a forecaster and a scheduling input and tests both on \DataSessions{} measurement sessions of a Starlink, a OneWeb, and a 5G link on one testbed~\cite{islam2025design}. Its contributions are the following.
\begin{enumerate}
  \item \textbf{What is predictable.} We separate the level, the change, and the spread of sub-second RTT. Past changes do not predict the next change beyond 0.4\,s, the RTT returns toward its recent level, and the spread is a property of the link state that the recent input reveals. Errors are skewed upward. We show that, as a result, intervals of constant width over-cover quiet windows and under-cover volatile ones on all six links.
  \item \textbf{A causality audit.} Aligning ping streams by time interpolation leaks part of the next sample into the input. We prove its effect on persistence and measure that a forecaster trained and tested on such data reaches a 0.2\,s error \LeakOptDSMAEhOne\% higher than reported once it is deployed.
  \item \textbf{A calibrated distributional forecaster.} A recurrent encoder-decoder anchored on the last observation outputs a mean and five non-crossing quantiles per step and target. Conformalized quantile regression (CQR)~\cite{romano2019conformalized} calibrates the quantiles on sessions held out from training.
  \item \textbf{An evaluation that bounds its own claims.} Against the same network with a scalar output, the 95\% interval score improves on all six links by \BSISWorstAbs\% to \BSISBestAbs\%, with session-bootstrap confidence intervals that exclude zero. We compare with the strongest alternatives we could build, namely asymmetric and volatility-scaled conformal intervals, Gaussian output heads, and a linear quantile model, and show how much of the gain each one recovers. In an emulated deadline-aware duplication scheduler with all settings fixed on held-out sessions, the calibrated bound needs \SchedSaveSLDL\% less duplication on Starlink DL with no more late packets, ties on two links, and loses on one.
\end{enumerate}

Section~\ref{sec:related} reviews related work and Section~\ref{sec:system} describes the system and the data. Section~\ref{sec:predictable} analyzes what is predictable and Section~\ref{sec:causal} the causality of the data pipeline. Section~\ref{sec:method} presents the forecaster, Sections~\ref{sec:setup} and~\ref{sec:results} the evaluation, and Section~\ref{sec:discussion} concludes the paper.
